\documentclass[11pt,a4paper]{article}

% ---------------------------------------------------------------
%  Research Statement -- Stanisław Szufa
%  EPFL IC Faculty Search 2026
% ---------------------------------------------------------------
\usepackage[T1]{fontenc}
\usepackage[utf8]{inputenc}
\usepackage{mathpazo}
\renewcommand{\ttdefault}{lmtt}  % outline typewriter font (avoids bitmap ectt)
\usepackage[hmargin=2.5cm,vmargin=2.2cm]{geometry}
\usepackage{microtype}
\usepackage{enumitem}
\usepackage{titlesec}
\usepackage{xcolor}
\usepackage[hidelinks]{hyperref}

\definecolor{accent}{RGB}{160,20,30}   % EPFL-ish red; change freely
\titleformat{\section}{\large\bfseries\color{accent}}{}{0pt}{}
\titlespacing*{\section}{0pt}{9pt}{3pt}
\titleformat{\subsection}[runin]{\normalsize\bfseries}{}{0pt}{}[.]
\titlespacing*{\subsection}{0pt}{5pt}{0.6em}
\setlength{\parindent}{0pt}
\setlength{\parskip}{3pt}
\setlist{nosep,leftmargin=1.2em}

\newcommand{\direction}[1]{\textbf{\color{accent}#1}}
% Reference to my own papers, using the keys of the publication list ([C3], [J4], ...).
\newcommand{\pub}[1]{~[#1]}
% Open problem inside a research direction.
\newcommand{\problem}[1]{\item \textbf{#1}}

\begin{document}

\begin{center}
  {\LARGE\bfseries Research Statement}\\[3pt]
  {\large Stanisław Szufa}
\end{center}
\vspace{2pt}

% ===============================================================
\section{Overview: Measuring Representation in Collective Decisions}
% ===============================================================

Many collective decisions are now made with the help of algorithms: a voting
rule decides which projects a city funds from its participatory budget,
electoral systems turn votes into parliaments, and large language models are
being tested for summarizing public consultations and drafting statements
that a whole group can agree on. In each case, the question is whether the
outcome represents the people it affects. Answering it with evidence
requires a way to compare decision problems (are synthetic test cases similar
to real ones?) and real data to test against. Computational social choice,
the field that studies voting rules with the tools of computer science, long
relied on proofs and on experiments with synthetic votes, often drawn
uniformly at random, with no general method to check how close these were to
real elections.

My work provides both. It falls into three lines, which also structure this
statement:
\begin{itemize}
  \item \emph{A geometry of collective decisions.} I introduced \emph{maps of
  elections}, which place elections in a common space, using a distance that I
  proposed, so that similar elections lie close together and one can see
  which models of synthetic votes resemble real data\pub{C3, J4}.
  \item \emph{Real data and real deployments.} I initiated Pabulib, the
  largest open repository of participatory budgeting data. It holds over
  2{,}000 elections from 8 countries and has been cited in more than 100
  papers. I am also a member of the team that convinced Polish cities to adopt
  the Method of Equal Shares, a proportional rule for participatory budgets.
  \item \emph{Toward language.} I now study collective decisions in which text
  and language models play a role, including simulated deliberation among LLM
  agents\pub{C28} and the voting record of the UN General Assembly\pub{P2}.
\end{itemize}
As a PhD student, I was the principal investigator of a National Science
Centre grant (NCN Preludium) on the \emph{election isomorphism problem}, from
which maps of elections grew. My thesis, \emph{Map of Elections} (advisor:
Piotr Faliszewski), received the Polish Prime Minister's Award (2025); I also
received a START fellowship (2023) and the AAMAS 2023 Best Student Paper
Award. I was then a CNRS postdoc at LAMSADE with J\'er\^ome Lang (2024--2026),
and I am now a postdoc at the University of Geneva.

\textbf{At EPFL, I will bring these tools (distances between decision
problems, real data, and large-scale experiments) to collective decisions in
which opinions are written in text and aggregated by AI}, such as AI-written
summaries of public consultations and consensus statements. We can already
generate such texts, but we lack principled ways to check whether they
represent the people behind them. For example, given the comments that
residents submitted in a consultation and an AI-written summary of them, how
can we measure which groups the summary leaves out, with guarantees, and test
the measure on real data?

% ===============================================================
\section{Past and Current Research}
% ===============================================================

\subsection{A geometry of collective decisions}
How similar are two elections? A useful distance between elections must
ignore what the candidates are called and in which
order the voters are listed, since two elections can have the same structure
with entirely different candidates (presidential candidates vs.\ types of
sushi). We showed that the most natural such distances, and the related
problem of finding one election inside another, are computationally
hard\pub{C2, J3, C5, J2}. I proposed instead the \emph{positionwise distance},
which compares two elections by how often each candidate is ranked in each
position, after optimally matching the candidates of the two elections. It
can be computed efficiently, and it made it possible to draw the \emph{map of
elections}\pub{C3}: compute the distances between all pairs in a collection
of elections and draw them in the plane, so that similar elections lie close
together. We then added landmark elections that make its regions
interpretable\pub{C4}; the complete framework appeared in \emph{Artificial
Intelligence}\pub{J4}.

Maps also show which parts of the space an experiment covers, which it
misses, and where real elections lie relative to synthetic ones. For example, the map showed that real-life
elections occupy only a small region of the space and that elections from the
Mallows model pass through it, so suitable Mallows parameters generate
realistic synthetic data. It also showed that the standard Mallows model drifts
toward identical votes as the number of candidates grows, which led us to
normalize its parameter\pub{C4}, and that ILP solvers for some multiwinner
rules take longest on elections close to impartial culture\pub{T1}.
The same approach works beyond ordinal elections. We built maps of approval
elections\pub{C7}, of elections of different sizes\pub{C19}, of stable
roommates and stable matching instances\pub{C9, J1} (Best Student Paper
Award), of tournaments\pub{C22}, and of fair division instances\pub{C25},
and used the geometry for preference learning\pub{C8}. To help others use
this methodology, I co-authored a guide to numerical experiments on
elections\pub{C15} and released two Python libraries: \texttt{mapof}, for
drawing maps, and \texttt{prefsampling}, for sampling synthetic preferences.

I also measure the structure within a single election, and these measures are
the starting point for the auditing tools I propose below. We introduced indices of agreement, diversity, and polarization based on
\emph{$k$-Kemeny}, which seeks $k$ representative rankings instead of one
consensus\pub{C10, C27, C24, C26}. We also look for
hidden structure in preference data, such as consistent
subelections\pub{C13} or the most conflicting pair of
candidates\pub{C18}; finding the groups of voters that an outcome serves
badly is a question of the same kind.

\subsection{Real data and real deployments}
In participatory budgeting (PB), residents vote on how to spend part of a
city's budget. I initiated Pabulib in 2020, and I built and
maintain it\pub{C11}. Our
later work on PB includes learning realistic models of approval
voting\pub{C20}, measuring how well individual projects performed under
a given rule\pub{C17}, and a game-theoretic analysis of how proposers
set project costs under the Method of Equal Shares\pub{C23}.
Under this rule, each resident
receives an equal share of the budget, and a project is funded only if its
supporters can jointly pay for it from their shares, so that each group of
residents gets a fair part of the budget. The team I belong to convinced
Polish cities, including Wieliczka and \'Swiecie, to adopt it.

\subsection{Toward language}
At the University of Geneva, I study collective decisions
in international politics, combining voting records with NLP pipelines that
extract structured information from large collections of diplomatic and policy
texts.

\emph{Simulating democratic deliberation with LLM agents.} I am a coauthor of
\emph{VoteSim}\pub{C28}, developed with collaborators at the University of
Toronto and Google DeepMind, which simulates a full democratic cycle with LLM
voters based on real participants' profiles. Across seven language models,
the voters preferred the bills of parliaments elected under proportional
systems to those of majoritarian ones.

\emph{What an abstention means in the UN General Assembly.} Abstention is the
most common way states withhold support in the General Assembly (11.8\% of
roughly 860{,}000 votes in 1946--2025), yet standard measures of state
preferences treat every abstention alike. In ongoing work that I
lead\pub{P2}, we classify all 101{,}070 abstentions by the context in which
they were cast, and show that this context tracks international conflict while
the overall abstention rate does not. We label each resolution's subject with a
language model, so the record becomes a large collection of votes over texts,
the kind of data that my agenda needs.

% ===============================================================
\section{Research Agenda}
% ===============================================================

Social choice models treat the alternatives as abstract symbols, whereas a
city project, a comment in an online consultation, or a UN resolution is a
text with content. Fish et al.\ (2024) prove proportionality guarantees
for slates of LLM-generated statements, assuming that the LLM answers its
queries accurately, and Tessler et al.\ (2024, in \emph{Science}) found that,
in groups of five, statements by an LLM mediator were accepted more often than
those of human mediators. In a position paper from a Dagstuhl seminar, we chart
how text can enter every stage of collective decision making\pub{P1}.
What is still missing are ways to \emph{evaluate} text-based outcomes and
\emph{real data} to evaluate them on. My agenda has three connected research directions, each continuing one of
the three lines of work above.

\direction{Direction A: A geometry of collective decision problems, from ballots to text.}
The goal is to extend the geometry of my earlier work to much larger instances
and to texts, so that experiments on both can be designed and compared
systematically.
\begin{itemize}
  \problem{Computing invariant distances at scale.} The natural
  isomorphism-invariant distances are hard to compute, and the positionwise
  distance is an efficient substitute for ordinal elections. Its cost,
  however, grows quickly with the number of candidates, and a map needs the
  distance between every pair of elections. Can sketching and sublinear
  techniques make maps feasible for elections with many candidates, such as
  large participatory budgets, and for very large collections of elections?
  For which other kinds of instances do efficient substitutes exist that give
  nearly the same maps as the exact distances?
  \problem{Benchmarks with provable coverage.} How many instances, and which
  ones, must a benchmark contain to cover the space of instances up to a given
  precision under a given distance? Which samplers achieve this? The answers
  will become defaults in \texttt{prefsampling} and \texttt{mapof}.
  \problem{Maps of deliberation.} When alternatives are texts, two decision
  problems can be similar in how people voted, in what the alternatives say,
  or both. One approach is to represent each alternative both by its votes
  and by an embedding of its text, and to match the alternatives of two
  instances with a cost that combines the two. Such distances will make it
  possible to compare
  deliberations, human or AI, the way we now compare elections, and to place
  sessions of the UN General Assembly on a common map. Direction~C will use
  these distances to compare text-based outcomes.
\end{itemize}

\direction{Direction B: Participatory democracy as a living laboratory.}
The goal is to turn real participatory processes into data for research on
text-based decisions, and to test in practice what theory predicts about
proportional rules.
\begin{itemize}
  \problem{Votes over texts.} I will extend Pabulib with project
  descriptions, residents' comments, and outcomes over several years, turning
  it into a longitudinal benchmark in which every ballot refers to a text.
  The UN General Assembly record, with 5{,}693 resolutions adopted by
  recorded vote together with their full texts\pub{P2}, is a
  complementary testbed: not a participatory process, but a large collection
  of votes over texts. Direction~C will use both as real data.
  \problem{Does proportionality carry over to content?} Proportional rules
  for participatory budgeting have been observed to spread funds evenly
  across project categories (Nelissen, 2023; Peters and Skowron, 2025). With project texts, we can
  ask whether they also give each group of residents a proportional share of
  the \emph{topics} it cares about, and design rules that guarantee it.
  \problem{Does the rule change the process?} We have analyzed how proposers
  choose project costs under the Method of Equal Shares\pub{C23}.
  Field studies with partner cities can test how the rule affects real
  proposers and participants: who proposes, what they propose, and who votes.
\end{itemize}
I will seek partnerships with Swiss municipalities for the field studies,
building on my experience of working with Polish cities on the Method of
Equal Shares. Until such partnerships are in place, the studies can start with
Polish cities that use the Method of Equal Shares and with existing Pabulib
data.

\direction{Direction C: Representation when opinions are expressed in text and aggregated by AI.}
The goal is to make it possible to check, with guarantees and on real data,
whether text aggregated by AI represents the people behind it.
\begin{itemize}
  \problem{Auditing summaries and consensus statements.} Given opinions
  written in natural language and a summary or consensus statement, how
  should we measure how proportionally it represents the population? Which
  measures satisfy natural axioms, and which can be estimated with few LLM
  queries and provable error bounds? A natural starting point is to
  treat the aggregate as a set of statements and to estimate, with LLM
  queries, how strongly each participant would endorse each of them, as in
  generative social choice; proportionality axioms for committees then become
  measurable properties of the text, which real consultation data can test.
  Because the endorsements are themselves estimated by a language model, the
  measures must be validated against participants' own judgments wherever
  these are available.
  \problem{An LLM aggregator is a voting rule.} Fix a language model and a
  prompt that turns a set of opinions into $k$ statements, and we obtain a
  multi-winner voting rule whose input is text. Existing methods, such as
  generative social choice and LLM mediators like the Habermas Machine, can
  then be tested like any other rule: does the outcome change when the same
  opinions are presented in a different order (anonymity)? Can a participant
  steer it by writing a manipulative opinion (manipulability)? Does it cover
  each group in proportion to its size? I helped establish this kind of
  large-scale numerical experiment for classical rules\pub{C15}. The result
  will be a public test suite that any proposed aggregator can be run against.
  \problem{Filling in sparse votes with text.} In online consultations, each
  participant reacts to only a small fraction of the comments; on Pol.is,
  about 9\%\pub{P1}. Language models can predict the missing reactions from
  the texts, but their errors then propagate to the outcome. How large is this
  effect, and which aggregation rules are robust to it?
\end{itemize}

\textbf{Goal for the first five years.} By the end of this period, my group
will have produced three things: maps that let researchers compare
deliberations the way they now compare elections (A); a public benchmark of
votes over texts, built from participatory budgets and the UN General Assembly
record (B); and open tools, with known guarantees, for auditing whether
AI-written summaries and consensus statements represent the people they
summarize (C).
These are the same kinds of contributions that my earlier work made for
ballots: data, measures, and an experimental method.

\textbf{Students, funding, and service.} Possible first PhD topics are scalable invariant distances (A, theoretical), proportional
coverage of topics in participatory budgeting (B, applied and data-centric),
and testing LLM aggregators as voting rules (C, empirical and NLP-oriented). I will apply to the Swiss National Science Foundation early on and for an ERC
Starting Grant built around Direction~C. I serve on the editorial board of
JAIR (since 2026) and on the program committees of NeurIPS, AAAI, AAMAS,
and IJCAI.

% ===============================================================
\section{Why EPFL}
% ===============================================================

EPFL does not have a group in computational social choice, but it is strong
in algorithms, natural language processing, and secure systems, which my
agenda relies on. It connects in particular with the work of three
colleagues:
\begin{itemize}
  \item \emph{Michael Kapralov}'s sublinear algorithms and sketching, for
  computing maps of elections with many candidates and of very large
  collections of elections (Direction~A);
  \item \emph{Bryan Ford}'s work on proof of personhood and coercion-resistant
  e-voting, for digital participatory processes that count each resident once
  and protect the secrecy of the vote (Direction~B);
  \item \emph{Antoine Bosselut}'s NLP group and EPFL's work on open
  multilingual language models such as Apertus, for auditing whether
  AI-written summaries represent people who write in different languages,
  which matters in a country with four national languages (Direction~C).
\end{itemize}
I also want to do this work in Switzerland because of its strong tradition of
voting. Swiss citizens decide federal, cantonal, and municipal questions by
popular vote several times a year, and in 2023 the city of Aarau used the
Method of Equal Shares for its participatory budget. Geneva, where I already
work, hosts the UN Human Rights Council, whose votes could extend my work on
the General Assembly\pub{P2}.

{\footnotesize\textbf{Other references.}
Fish, G\"olz, Parkes, Procaccia, Rusak, Shapira, and W\"uthrich. Generative
social choice. \emph{Proceedings of EC-2024}.
Nelissen. An empirical analysis of participatory budgeting in Amsterdam.
arXiv:2310.18033, 2023.
Peters and Skowron. The Method of Equal Shares for participatory budgeting.
\url{https://equalshares.net}, 2025.
Tessler, Bakker, Jarrett, et al. AI can help humans find common ground in
democratic deliberation. \emph{Science} 386(6719), 2024.\par}

\end{document}
